\documentclass[12pt, a4paper]{article}

\usepackage[a4paper, margin=1in]{geometry}

\usepackage[utf8]{inputenc}
\usepackage[mathscr]{euscript}
\usepackage{amssymb,amsmath,amsthm,amsfonts}
\usepackage[shortlabels]{enumitem}
\usepackage{multicol,multirow}
\usepackage{lipsum}
\usepackage{balance}
\usepackage{calc}
\usepackage[colorlinks=true,citecolor=blue,linkcolor=blue]{hyperref}
\usepackage{import}
\usepackage{xifthen}
\usepackage{pdfpages}
\usepackage{transparent}
\usepackage{listings}

\newcommand{\incfig}[2][1.0]{
    \def\svgwidth{#1\columnwidth}
    \import{./figures/}{#2.pdf_tex}
}

\newcommand{\incimg}[2][1.0]{
  \includegraphics[width=#1\columnwidth]{./figures/#2}
}

\newlist{enumproof}{enumerate}{4}
\setlist[enumproof,1]{label=\arabic*., parsep=1em}
\setlist[enumproof,2]{label=\arabic{enumproofi}.\arabic*., parsep=1em}
\setlist[enumproof,3]{label=\arabic{enumproofi}.\arabic{enumproofii}.\arabic*., parsep=1em}
\setlist[enumproof,4]{label=\arabic{enumproofi}.\arabic{enumproofii}.\arabic{enumproofiii}.\arabic*., parsep=1em}

\renewcommand{\qedsymbol}{\ensuremath{\blacksquare}}

\lstdefinestyle{mystyle}{
  language=Python, % Set the language to Python
  commentstyle=\color{codegray}\itshape, % Python comments (#)
  keywordstyle=\color{orange}\bfseries,  % Python reserved keywords
  stringstyle=\color{mauve},             % Strings
  basicstyle={\ttfamily\footnotesize},   % Monospace font
  breakatwhitespace=true,         
  breaklines=true,                 
  captionpos=b,                    
  keepspaces=true,                 
  numbers=none,                    
  tabsize=2,
  showstringspaces=false,
  % Extra keyword groups for Python
  morekeywords=[2]{self, cls, __init__, __str__, __repr__}, 
  keywordstyle=[2]\color{identblue},     % Special identifiers / methods
  morekeywords=[3]{int, float, bool, str, list, dict, set, tuple, object, None, True, False},
  keywordstyle=[3]\color{codegreen},     % Built-in types
  morekeywords=[4]{print, input, open, range, len, type, isinstance, super, enumerate, zip, map, filter, sum, min, max, sorted},
  keywordstyle=[4]\color{cyan},          % Built-in functions
}

% Things Lie
\newcommand{\kb}{\mathfrak b}
\newcommand{\kg}{\mathfrak g}
\newcommand{\kh}{\mathfrak h}
\newcommand{\kn}{\mathfrak n}
\newcommand{\ku}{\mathfrak u}
\newcommand{\kz}{\mathfrak z}
\DeclareMathOperator{\Ext}{Ext} % Ext functor
\DeclareMathOperator{\Tor}{Tor} % Tor functor
\newcommand{\gl}{\opname{\mathfrak{gl}}} % frak gl group
\renewcommand{\sl}{\opname{\mathfrak{sl}}} % frak sl group chktex 6

% More script letters etc.
\newcommand{\SA}{\mathcal A}
\newcommand{\SB}{\mathcal B}
\newcommand{\SC}{\mathcal C}
\newcommand{\SF}{\mathcal F}
\newcommand{\SG}{\mathcal G}
\newcommand{\SH}{\mathcal H}
\newcommand{\OO}{\mathcal O}

\newcommand{\SCA}{\mathscr A}
\newcommand{\SCB}{\mathscr B}
\newcommand{\SCC}{\mathscr C}
\newcommand{\SCD}{\mathscr D}
\newcommand{\SCE}{\mathscr E}
\newcommand{\SCF}{\mathscr F}
\newcommand{\SCG}{\mathscr G}
\newcommand{\SCH}{\mathscr H}

% Mathfrak primes
\newcommand{\km}{\mathfrak m}
\newcommand{\kp}{\mathfrak p}
\newcommand{\kq}{\mathfrak q}

% number sets
\newcommand{\RR}[1][]{\ensuremath{\ifstrempty{#1}{\mathbb{R}}{\mathbb{R}^{#1}}}}
\newcommand{\NN}[1][]{\ensuremath{\ifstrempty{#1}{\mathbb{N}}{\mathbb{N}^{#1}}}}
\newcommand{\ZZ}[1][]{\ensuremath{\ifstrempty{#1}{\mathbb{Z}}{\mathbb{Z}^{#1}}}}
\newcommand{\QQ}[1][]{\ensuremath{\ifstrempty{#1}{\mathbb{Q}}{\mathbb{Q}^{#1}}}}
\newcommand{\CC}[1][]{\ensuremath{\ifstrempty{#1}{\mathbb{C}}{\mathbb{C}^{#1}}}}
\newcommand{\PP}[1][]{\ensuremath{\ifstrempty{#1}{\mathbb{P}}{\mathbb{P}^{#1}}}}
\newcommand{\HH}[1][]{\ensuremath{\ifstrempty{#1}{\mathbb{H}}{\mathbb{H}^{#1}}}}
\newcommand{\FF}[1][]{\ensuremath{\ifstrempty{#1}{\mathbb{F}}{\mathbb{F}^{#1}}}}
% expected value
\newcommand{\EE}{\ensuremath{\mathbb{E}}}
\newcommand{\charin}{\text{ char }}
\DeclareMathOperator{\sign}{sign}
\DeclareMathOperator{\Aut}{Aut}
\DeclareMathOperator{\Inn}{Inn}
\DeclareMathOperator{\Syl}{Syl}
\DeclareMathOperator{\Gal}{Gal}
\DeclareMathOperator{\GL}{GL} % General linear group
\DeclareMathOperator{\SL}{SL} % Special linear group

%---------------------------------------
% BlackBoard Math Fonts :-
%---------------------------------------

%Captital Letters
\newcommand{\bbA}{\mathbb{A}}	\newcommand{\bbB}{\mathbb{B}}
\newcommand{\bbC}{\mathbb{C}}	\newcommand{\bbD}{\mathbb{D}}
\newcommand{\bbE}{\mathbb{E}}	\newcommand{\bbF}{\mathbb{F}}
\newcommand{\bbG}{\mathbb{G}}	\newcommand{\bbH}{\mathbb{H}}
\newcommand{\bbI}{\mathbb{I}}	\newcommand{\bbJ}{\mathbb{J}}
\newcommand{\bbK}{\mathbb{K}}	\newcommand{\bbL}{\mathbb{L}}
\newcommand{\bbM}{\mathbb{M}}	\newcommand{\bbN}{\mathbb{N}}
\newcommand{\bbO}{\mathbb{O}}	\newcommand{\bbP}{\mathbb{P}}
\newcommand{\bbQ}{\mathbb{Q}}	\newcommand{\bbR}{\mathbb{R}}
\newcommand{\bbS}{\mathbb{S}}	\newcommand{\bbT}{\mathbb{T}}
\newcommand{\bbU}{\mathbb{U}}	\newcommand{\bbV}{\mathbb{V}}
\newcommand{\bbW}{\mathbb{W}}	\newcommand{\bbX}{\mathbb{X}}
\newcommand{\bbY}{\mathbb{Y}}	\newcommand{\bbZ}{\mathbb{Z}}

%---------------------------------------
% MathCal Fonts :-
%---------------------------------------

%Captital Letters
\newcommand{\mcA}{\mathcal{A}}	\newcommand{\mcB}{\mathcal{B}}
\newcommand{\mcC}{\mathcal{C}}	\newcommand{\mcD}{\mathcal{D}}
\newcommand{\mcE}{\mathcal{E}}	\newcommand{\mcF}{\mathcal{F}}
\newcommand{\mcG}{\mathcal{G}}	\newcommand{\mcH}{\mathcal{H}}
\newcommand{\mcI}{\mathcal{I}}	\newcommand{\mcJ}{\mathcal{J}}
\newcommand{\mcK}{\mathcal{K}}	\newcommand{\mcL}{\mathcal{L}}
\newcommand{\mcM}{\mathcal{M}}	\newcommand{\mcN}{\mathcal{N}}
\newcommand{\mcO}{\mathcal{O}}	\newcommand{\mcP}{\mathcal{P}}
\newcommand{\mcQ}{\mathcal{Q}}	\newcommand{\mcR}{\mathcal{R}}
\newcommand{\mcS}{\mathcal{S}}	\newcommand{\mcT}{\mathcal{T}}
\newcommand{\mcU}{\mathcal{U}}	\newcommand{\mcV}{\mathcal{V}}
\newcommand{\mcW}{\mathcal{W}}	\newcommand{\mcX}{\mathcal{X}}
\newcommand{\mcY}{\mathcal{Y}}	\newcommand{\mcZ}{\mathcal{Z}}


%---------------------------------------
% Bold Math Fonts :-
%---------------------------------------

%Captital Letters
\newcommand{\bmA}{\boldsymbol{A}}	\newcommand{\bmB}{\boldsymbol{B}}
\newcommand{\bmC}{\boldsymbol{C}}	\newcommand{\bmD}{\boldsymbol{D}}
\newcommand{\bmE}{\boldsymbol{E}}	\newcommand{\bmF}{\boldsymbol{F}}
\newcommand{\bmG}{\boldsymbol{G}}	\newcommand{\bmH}{\boldsymbol{H}}
\newcommand{\bmI}{\boldsymbol{I}}	\newcommand{\bmJ}{\boldsymbol{J}}
\newcommand{\bmK}{\boldsymbol{K}}	\newcommand{\bmL}{\boldsymbol{L}}
\newcommand{\bmM}{\boldsymbol{M}}	\newcommand{\bmN}{\boldsymbol{N}}
\newcommand{\bmO}{\boldsymbol{O}}	\newcommand{\bmP}{\boldsymbol{P}}
\newcommand{\bmQ}{\boldsymbol{Q}}	\newcommand{\bmR}{\boldsymbol{R}}
\newcommand{\bmS}{\boldsymbol{S}}	\newcommand{\bmT}{\boldsymbol{T}}
\newcommand{\bmU}{\boldsymbol{U}}	\newcommand{\bmV}{\boldsymbol{V}}
\newcommand{\bmW}{\boldsymbol{W}}	\newcommand{\bmX}{\boldsymbol{X}}
\newcommand{\bmY}{\boldsymbol{Y}}	\newcommand{\bmZ}{\boldsymbol{Z}}
%Small Letters
\newcommand{\bma}{\boldsymbol{a}}	\newcommand{\bmb}{\boldsymbol{b}}
\newcommand{\bmc}{\boldsymbol{c}}	\newcommand{\bmd}{\boldsymbol{d}}
\newcommand{\bme}{\boldsymbol{e}}	\newcommand{\bmf}{\boldsymbol{f}}
\newcommand{\bmg}{\boldsymbol{g}}	\newcommand{\bmh}{\boldsymbol{h}}
\newcommand{\bmi}{\boldsymbol{i}}	\newcommand{\bmj}{\boldsymbol{j}}
\newcommand{\bmk}{\boldsymbol{k}}	\newcommand{\bml}{\boldsymbol{l}}
\newcommand{\bmm}{\boldsymbol{m}}	\newcommand{\bmn}{\boldsymbol{n}}
\newcommand{\bmo}{\boldsymbol{o}}	\newcommand{\bmp}{\boldsymbol{p}}
\newcommand{\bmq}{\boldsymbol{q}}	\newcommand{\bmr}{\boldsymbol{r}}
\newcommand{\bms}{\boldsymbol{s}}	\newcommand{\bmt}{\boldsymbol{t}}
\newcommand{\bmu}{\boldsymbol{u}}	\newcommand{\bmv}{\boldsymbol{v}}
\newcommand{\bmw}{\boldsymbol{w}}	\newcommand{\bmx}{\boldsymbol{x}}
\newcommand{\bmy}{\boldsymbol{y}}	\newcommand{\bmz}{\boldsymbol{z}}

%---------------------------------------
% Scr Math Fonts :-
%---------------------------------------

\newcommand{\sA}{{\mathscr{A}}}   \newcommand{\sB}{{\mathscr{B}}}
\newcommand{\sC}{{\mathscr{C}}}   \newcommand{\sD}{{\mathscr{D}}}
\newcommand{\sE}{{\mathscr{E}}}   \newcommand{\sF}{{\mathscr{F}}}
\newcommand{\sG}{{\mathscr{G}}}   \newcommand{\sH}{{\mathscr{H}}}
\newcommand{\sI}{{\mathscr{I}}}   \newcommand{\sJ}{{\mathscr{J}}}
\newcommand{\sK}{{\mathscr{K}}}   \newcommand{\sL}{{\mathscr{L}}}
\newcommand{\sM}{{\mathscr{M}}}   \newcommand{\sN}{{\mathscr{N}}}
\newcommand{\sO}{{\mathscr{O}}}   \newcommand{\sP}{{\mathscr{P}}}
\newcommand{\sQ}{{\mathscr{Q}}}   \newcommand{\sR}{{\mathscr{R}}}
\newcommand{\sS}{{\mathscr{S}}}   \newcommand{\sT}{{\mathscr{T}}}
\newcommand{\sU}{{\mathscr{U}}}   \newcommand{\sV}{{\mathscr{V}}}
\newcommand{\sW}{{\mathscr{W}}}   \newcommand{\sX}{{\mathscr{X}}}
\newcommand{\sY}{{\mathscr{Y}}}   \newcommand{\sZ}{{\mathscr{Z}}}


%---------------------------------------
% Math Fraktur Font
%---------------------------------------

%Captital Letters
\newcommand{\mfA}{\mathfrak{A}}	\newcommand{\mfB}{\mathfrak{B}}
\newcommand{\mfC}{\mathfrak{C}}	\newcommand{\mfD}{\mathfrak{D}}
\newcommand{\mfE}{\mathfrak{E}}	\newcommand{\mfF}{\mathfrak{F}}
\newcommand{\mfG}{\mathfrak{G}}	\newcommand{\mfH}{\mathfrak{H}}
\newcommand{\mfI}{\mathfrak{I}}	\newcommand{\mfJ}{\mathfrak{J}}
\newcommand{\mfK}{\mathfrak{K}}	\newcommand{\mfL}{\mathfrak{L}}
\newcommand{\mfM}{\mathfrak{M}}	\newcommand{\mfN}{\mathfrak{N}}
\newcommand{\mfO}{\mathfrak{O}}	\newcommand{\mfP}{\mathfrak{P}}
\newcommand{\mfQ}{\mathfrak{Q}}	\newcommand{\mfR}{\mathfrak{R}}
\newcommand{\mfS}{\mathfrak{S}}	\newcommand{\mfT}{\mathfrak{T}}
\newcommand{\mfU}{\mathfrak{U}}	\newcommand{\mfV}{\mathfrak{V}}
\newcommand{\mfW}{\mathfrak{W}}	\newcommand{\mfX}{\mathfrak{X}}
\newcommand{\mfY}{\mathfrak{Y}}	\newcommand{\mfZ}{\mathfrak{Z}}
%Small Letters
\newcommand{\mfa}{\mathfrak{a}}	\newcommand{\mfb}{\mathfrak{b}}
\newcommand{\mfc}{\mathfrak{c}}	\newcommand{\mfd}{\mathfrak{d}}
\newcommand{\mfe}{\mathfrak{e}}	\newcommand{\mff}{\mathfrak{f}}
\newcommand{\mfg}{\mathfrak{g}}	\newcommand{\mfh}{\mathfrak{h}}
\newcommand{\mfi}{\mathfrak{i}}	\newcommand{\mfj}{\mathfrak{j}}
\newcommand{\mfk}{\mathfrak{k}}	\newcommand{\mfl}{\mathfrak{l}}
\newcommand{\mfm}{\mathfrak{m}}	\newcommand{\mfn}{\mathfrak{n}}
\newcommand{\mfo}{\mathfrak{o}}	\newcommand{\mfp}{\mathfrak{p}}
\newcommand{\mfq}{\mathfrak{q}}	\newcommand{\mfr}{\mathfrak{r}}
\newcommand{\mfs}{\mathfrak{s}}	\newcommand{\mft}{\mathfrak{t}}
\newcommand{\mfu}{\mathfrak{u}}	\newcommand{\mfv}{\mathfrak{v}}
\newcommand{\mfw}{\mathfrak{w}}	\newcommand{\mfx}{\mathfrak{x}}
\newcommand{\mfy}{\mathfrak{y}}	\newcommand{\mfz}{\mathfrak{z}}


\newcommand{\mytitle}{FMA\_3TC13\_TP Poly Exercices 5}
\newcommand{\myauthor}{github/omgeta}
\newcommand{\mydate}{AU 25-26 S1}

\begin{document}
\raggedright
\footnotesize
\begin{center}
{\normalsize{\textbf{\mytitle}}} \\
{\footnotesize{\mydate\hspace{2pt}\textemdash\hspace{2pt}\myauthor}}
\end{center}
\setlist{topsep=-1em, itemsep=-1em, parsep=2em}
\setlist[enumerate]{topsep=0.4em,itemsep=0.8em,parsep=0.35em}
%%%%%%%%%%%%%%%%%%%%%%%%%%%%%%%%%%%%%%%%%%%%%%%%%%%%%%
%                      Begin                         %
%%%%%%%%%%%%%%%%%%%%%%%%%%%%%%%%%%%%%%%%%%%%%%%%%%%%%%
\begin{enumerate}[label=5.\arabic*., start=54]

  \item \begin{minipage}[t]{\linewidth}\raggedright
\textbf{Énoncé.} Calculer les déterminants suivants.

\begin{enumerate}[\arabic*.]
        \item \textbf{Question.} $\displaystyle\det\begin{pmatrix}0&6&10&180\\6&14&12&215\\-9&-18&-14&-250\\3&0&-5&-90\end{pmatrix}$.
\par\smallskip
\textbf{Réponse.} Par élimination de Gauss (en tenant compte des échanges de lignes),

$\det A=30$.

  \item \textbf{Question.} $\displaystyle\det\begin{pmatrix}2&0&0&2\\1&1&-1&-3\\-1&0&1&4\\0&-1&0&1\end{pmatrix}$.
\par\smallskip
\textbf{Réponse.} En développant suivant la première ligne, $\det B=2\cdot2-2\cdot0=4$.

  \item \textbf{Question.} $\displaystyle\det\begin{pmatrix}-1&1&0&1\\-4&-1&-1&-2\\1&0&-1&-3\\-1&2&-3&1\end{pmatrix}$.
\par\smallskip
\textbf{Réponse.} $\det C=-44$.

\end{enumerate}
\end{minipage}

  \item \begin{minipage}[t]{\linewidth}\raggedright
\textbf{Question.} Soient $n=p+q$, $A\in K^{p\times p}$, $B\in K^{p\times q}$, $C\in K^{q\times q}$. Montrer que $\det\begin{pmatrix}A&B\\0&C\end{pmatrix}=\det A\det C$.
\par\smallskip
\textbf{Réponse.} Dans le développement du déterminant par permutations, un terme non nul doit associer les $q$ dernières lignes aux $q$ dernières colonnes, car le bloc inférieur gauche est nul.
    Les $p$ premières lignes sont donc associées aux $p$ premières colonnes.
    Les permutations et leurs signatures se séparent en deux blocs.

    $\therefore \det Z=\det A\det C$.
\end{minipage}

  \item \begin{minipage}[t]{\linewidth}\raggedright
\textbf{Question.} Soit $j=e^{2i\pi/3}$. Calculer $\displaystyle\det\begin{pmatrix}1&j&j^2\\j&1&j^2\\j^2&j&1\end{pmatrix}$.
\par\smallskip
\textbf{Réponse.} $1+j+j^2=0$. La somme des trois colonnes est nulle, donc les colonnes sont liées.
    $\therefore D=0$.
\end{minipage}

  \item \begin{minipage}[t]{\linewidth}\raggedright
\textbf{Question.} Pour quelles valeurs de $a\in\RR$ la matrice $\displaystyle A=\begin{pmatrix}a&-2&2\\1&-a&3\\5&-8&12a\end{pmatrix}$ est-elle inversible ?
\par\smallskip
\textbf{Réponse.} En développant,
    \[
      \det A=-2(a-1)(6a^2+6a-23)
    \]
    $\therefore A$ est inversible si et seulement si
    $\displaystyle a\notin\left\{1,\frac{-3-7\sqrt3}{6},\frac{-3+7\sqrt3}{6}\right\}$.
\end{minipage}

  \item \begin{minipage}[t]{\linewidth}\raggedright
\textbf{Énoncé.} Soient $a_1,\ldots,a_n\in\RR$.

\begin{enumerate}[\arabic*.]
        \item \textbf{Question.} Calculer le déterminant de Vandermonde $V_n=(a_j^{i-1})_{1\leq i,j\leq n}$.
\par\smallskip
\textbf{Réponse.} Comme polynôme en $a_n$, le déterminant s'annule en $a_1,\ldots,a_{n-1}$ et son coefficient dominant vaut $\det V_{n-1}$.
        Par récurrence,
        \[
          \therefore\det V_n=\prod_{1\leq i<j\leq n}(a_j-a_i)
        \]
        Si deux paramètres coïncident, les deux membres sont nuls.

  \item \textbf{Question.} En déduire $\displaystyle\det\begin{pmatrix}1&2a+3&3a^2+4a&4a^3+5a\\1&2b+3&3b^2+4b&4b^3+5b\\1&2c+3&3c^2+4c&4c^3+5c\\1&2d+3&3d^2+4d&4d^3+5d\end{pmatrix}$.
\par\smallskip
\textbf{Réponse.} Les colonnes sont des polynômes de degrés $0,1,2,3$, de coefficients dominants $1,2,3,4$.
        En retranchant les termes de degrés inférieurs,
        \[
          \therefore D=24(b-a)(c-a)(d-a)(c-b)(d-b)(d-c)
        \]

  \item \textbf{Question.} Calculer $\displaystyle\det\begin{pmatrix}b+c&a+c&a+b\\b^2+c^2&a^2+c^2&a^2+b^2\\b^3+c^3&a^3+c^3&a^3+b^3\end{pmatrix}$.
\par\smallskip
\textbf{Réponse.} La matrice s'écrit
        $\displaystyle\begin{pmatrix}a&b&c\\a^2&b^2&c^2\\a^3&b^3&c^3\end{pmatrix}
          \begin{pmatrix}0&1&1\\1&0&1\\1&1&0\end{pmatrix}$.
        $\therefore D=2abc(b-a)(c-a)(c-b)$.

  \item \textbf{Question.} Calculer $D=\det(\cos(j\theta_i))_{0\leq i,j\leq n}$, où la colonne $j=0$ est constante égale à $1$.
\par\smallskip
\textbf{Réponse.} $\cos(k\theta)=T_k(\cos\theta)$, où $T_0=1$, $T_1=X$ et $T_{k+1}=2XT_k-T_{k-1}$.
        Le coefficient dominant de $T_k$ est $2^{k-1}$ pour $k\geq1$.
        \[
          \therefore D=2^{n(n-1)/2}\prod_{0\leq i<j\leq n}(\cos\theta_j-\cos\theta_i)
        \]

\end{enumerate}
\end{minipage}

  \item \begin{minipage}[t]{\linewidth}\raggedright
\textbf{Énoncé.} Soient $A,B\in\RR^{n\times n}$ et $P\in\operatorname{GL}_n(\CC)$ telles que $B=P^{-1}AP$. On écrit $P=R+iS$ avec $R,S$ réelles.

\begin{enumerate}[\arabic*.]
        \item \begin{enumerate}[a.]
            \item \textbf{Question.} Vérifier $RB=AR$ et $SB=AS$.
\par\smallskip
\textbf{Réponse.} $PB=AP$ donne $(R+iS)B=AR+iAS$.
            En identifiant parties réelle et imaginaire, $RB=AR$ et $SB=AS$.

  \item \textbf{Question.} En déduire $(R+xS)B=A(R+xS)$ pour tout $x\in\CC$.
\par\smallskip
\textbf{Réponse.} Par combinaison linéaire, $(R+xS)B=A(R+xS)$ pour tout $x\in\CC$.

\end{enumerate}

  \item \textbf{Question.} Prouver qu’il existe $x_0\in\RR$ tel que $R+x_0S$ soit inversible.
\par\smallskip
\textbf{Réponse.} $p(X)=\det(R+XS)$ est un polynôme réel non nul, car $p(i)=\det P\neq0$.
        Il n'a qu'un nombre fini de racines, donc il existe $x_0\in\RR$ tel que $p(x_0)\neq0$.

  \item \textbf{Question.} Conclure que $A,B$ sont semblables sur $\RR$.
\par\smallskip
\textbf{Réponse.} Pour $Q=R+x_0S\in\operatorname{GL}_n(\RR)$, $QB=AQ$.
        $\therefore B=Q^{-1}AQ$ : $A$ et $B$ sont semblables sur $\RR$.

\end{enumerate}
\end{minipage}

  \item \begin{minipage}[t]{\linewidth}\raggedright
\textbf{Énoncé.} Soit $n\geq2$. On note $\operatorname{Com}(X)$ la comatrice de $X\in K^{n\times n}$.

\begin{enumerate}[\arabic*.]
        \item \textbf{Question.} Montrer que $\det(\operatorname{Com}(X))=(\det X)^{n-1}$.
\par\smallskip
\textbf{Réponse.} Si $X$ est inversible, $\operatorname{Com}(X)=\det(X)X^{-\top}$, d'où
        \[
          \det(\operatorname{Com}(X))=(\det X)^n/\det X=(\det X)^{n-1}
        \]
        L'identité reste vraie pour $X$ singulière : appliquer le calcul à $X+tI_n$ sur $K(t)$, puis évaluer l'identité polynomiale en $t=0$.

  \item \textbf{Question.} Pour $M\in\operatorname{GL}_n(K)$, résoudre $M=\operatorname{Com}(X)$.
\par\smallskip
\textbf{Réponse.} Si $M=\operatorname{Com}(X)$ est inversible, alors $t=\det X\neq0$ vérifie $t^{n-1}=\det M$ et $X=tM^{-\top}$.
        Réciproquement, chaque tel $t\in K$ fournit une solution.
        \[
          \therefore\mathcal S=\{tM^{-\top}:t\in K,\ t^{n-1}=\det M\}
        \]

\end{enumerate}
\end{minipage}

  \item \begin{minipage}[t]{\linewidth}\raggedright
\textbf{Énoncé.} Soient $j=-1/2+i\sqrt3/2$, $a,b,c\in\CC$, $\displaystyle A=\begin{pmatrix}a&b&c\\c&a&b\\b&c&a\end{pmatrix}$ et $\displaystyle J=\begin{pmatrix}1&1&1\\1&j&j^2\\1&j^2&j\end{pmatrix}$.

\begin{enumerate}[\arabic*.]
        \item \textbf{Question.} Calculer $JA$.
\par\smallskip
\textbf{Réponse.} Posons $s=a+b+c$, $p=a+jc+j^2b$ et $q=a+j^2c+jb$.

$\displaystyle JA=\begin{pmatrix}s&s&s\\p&jp&j^2p\\q&j^2q&jq\end{pmatrix}$.

  \item \textbf{Question.} Trouver une matrice diagonale $\Delta$ telle que $JA=\Delta J$.
\par\smallskip
\textbf{Réponse.} $JA=\Delta J$, avec $\Delta=\operatorname{diag}(s,p,q)$.

  \item \textbf{Question.} Donner une forme factorisée de $\det A$.
\par\smallskip
\textbf{Réponse.} $J$ est inversible (Vandermonde), donc
        \[
          \det A=spq=(a+b+c)(a+jb+j^2c)(a+j^2b+jc)
          =a^3+b^3+c^3-3abc
        \]

  \item \textbf{Question.} Pour $a,b,c\in\RR$, discuter et résoudre $ax+by+cz=1$, $cx+ay+bz=1$, $bx+cy+az=1$.
\par\smallskip
\textbf{Réponse.} Pour $a,b,c\in\RR$, $pq=a^2+b^2+c^2-ab-ac-bc
        =\big((a-b)^2+(a-c)^2+(b-c)^2\big)/2$.

        Si $s=0$, sommer les équations donne $0=3$, donc aucune solution.
        Si $s\neq0$ et $(a,b,c)$ ne sont pas tous égaux, l'unique solution est $(1/s,1/s,1/s)$.
        Si $a=b=c\neq0$, les solutions sont tous les triplets tels que $x+y+z=1/a$.

\end{enumerate}
\end{minipage}

  \item \begin{minipage}[t]{\linewidth}\raggedright
\textbf{Question.} Soient $a_0,\ldots,a_{n-1}\in K$. Calculer le déterminant de $xI_n-C$, où $C$ a des $1$ sur la sous-diagonale, la dernière colonne $(-a_0,\ldots,-a_{n-1})^\top$, et des $0$ ailleurs (matrice compagnon).
\par\smallskip
\textbf{Réponse.} En développant suivant la dernière colonne, le cofacteur de $a_k$ vaut $x^k$ : les signes des $-1$ et du cofacteur se compensent.
    Le terme diagonal supplémentaire donne $x^n$.
    \[
      \therefore\Delta_n=x^n+a_{n-1}x^{n-1}+\cdots+a_1x+a_0
    \]
\end{minipage}

  \item \begin{minipage}[t]{\linewidth}\raggedright
\textbf{Énoncé.} Soit $A=(a_{i,j})\in\RR^{n\times n}$ telle que $\sum_{i=1}^n|a_{i,j}|\leq1$ pour chaque colonne $j$.

\begin{enumerate}[\arabic*.]
        \item \textbf{Question.} Montrer que $|\det A|\leq1$.
\par\smallskip
\textbf{Réponse.} En notant $c_j$ les colonnes et en développant par multilinéarité,
        \[
          |\det A|\leq\sum_{\sigma\in S_n}\prod_j|a_{\sigma(j),j}|
          \leq\prod_{j=1}^n\left(\sum_i|a_{i,j}|\right)\leq1
        \]

  \item \textbf{Question.} Étudier le cas d’égalité.
\par\smallskip
\textbf{Réponse.} Si $|\det A|=1$, chaque colonne a une somme de valeurs absolues égale à $1$.
        Deux colonnes ne peuvent avoir un coefficient non nul sur une même ligne : sinon le produit des sommes contiendrait un terme strictement positif à indices répétés, absent du déterminant, et l'inégalité serait stricte.
        Les $n$ supports sont donc disjoints et non vides parmi $n$ lignes : chacun est un singleton.

        $\therefore A$ est une matrice de permutation dont les $1$ sont éventuellement remplacés par $-1$.
        Réciproquement, ces matrices ont un déterminant $\pm1$.

\end{enumerate}
\end{minipage}

  \item \begin{minipage}[t]{\linewidth}\raggedright
\textbf{Question.} Pour $a,b,c\in K$, calculer $\displaystyle D=\det\begin{pmatrix}0&1&1&1\\1&0&c^2&b^2\\1&c^2&0&a^2\\1&b^2&a^2&0\end{pmatrix}$.
\par\smallskip
\textbf{Réponse.} Après $C_3\leftarrow C_3-C_2$, $C_4\leftarrow C_4-C_2$, puis développement suivant la première ligne,
    \[
      D=(a^2-b^2-c^2)^2-4b^2c^2
       =(a-b-c)(a+b+c)(a-b+c)(a+b-c)
    \]
\end{minipage}

  \item \begin{minipage}[t]{\linewidth}\raggedright
\textbf{Question.} Pour $\theta\in\RR$ et $n\geq2$, calculer $\Delta=\det(\cos((i+j)\theta))_{0\leq i,j\leq n}$.
\par\smallskip
\textbf{Réponse.} Les trois premières colonnes vérifient $C_1+C_3=2\cos\theta\,C_2$.
    Elles sont liées. $\therefore\Delta=0$ pour $n\geq2$.
\end{minipage}

  \item \begin{minipage}[t]{\linewidth}\raggedright
\textbf{Énoncé.} Soient $a,b,c,d\in K$, $\displaystyle J=\begin{pmatrix}1&1&1&1\\1&1&-1&-1\\1&-1&-1&1\\1&-1&1&-1\end{pmatrix}$ et $\displaystyle A=\begin{pmatrix}a&b&c&d\\b&a&d&c\\c&d&a&b\\d&c&b&a\end{pmatrix}$.

\begin{enumerate}[\arabic*.]
        \item \textbf{Question.} Calculer $\det J$.
\par\smallskip
\textbf{Réponse.} Posons $s_1=a+b+c+d$, $s_2=a+b-c-d$, $s_3=a-b-c+d$, $s_4=a-b+c-d$.

Par élimination, $\det J=16$.

  \item \textbf{Question.} Calculer $AJ$.
\par\smallskip
\textbf{Réponse.} $AJ=J\operatorname{diag}(s_1,s_2,s_3,s_4)$.

  \item \textbf{Question.} En déduire $\det A$.
\par\smallskip
\textbf{Réponse.} $\det A=s_1s_2s_3s_4$.
        Si $\operatorname{car}K\neq2$, on prend les déterminants et on simplifie par $16$.
        En caractéristique $2$, le même résultat suit de l'identité polynomiale à coefficients entiers.

\end{enumerate}
\end{minipage}

  \item \begin{minipage}[t]{\linewidth}\raggedright
\textbf{Question.} Soient $a_i,b_i\in K$ avec $\prod_i a_i\neq0$. Calculer le déterminant de la matrice de coefficients $m_{i,j}=a_i\delta_{i,j}+b_i$.
\par\smallskip
\textbf{Réponse.} La $j$-ième colonne est $a_je_j+b$, où $b=(b_1,\ldots,b_n)^\top$.
    Par multilinéarité, les termes contenant deux colonnes $b$ s'annulent.
    Il reste le produit diagonal et les $n$ termes contenant une seule colonne $b$ :
    \[
      \therefore D=\prod_i a_i+\sum_i b_i\prod_{j\neq i}a_j
      =A\left(1+\sum_i\frac{b_i}{a_i}\right)
    \]
\end{minipage}

  \item \begin{minipage}[t]{\linewidth}\raggedright
\textbf{Question.} Pour $p,q\in K$, calculer $D_n=\det M$, où $m_{i,i}=p+q$, $m_{i,j}=p$ si $i<j$, $m_{i,j}=q$ si $i>j$.
\par\smallskip
\textbf{Réponse.} En effectuant $C_j\leftarrow C_j-C_{j-1}$ de droite à gauche, puis en développant, on obtient
    $D_n=(p+q)D_{n-1}-pqD_{n-2}$, avec $D_0=1$, $D_1=p+q$.
    Par récurrence,
    \[
      \therefore D_n=\sum_{k=0}^n p^{n-k}q^k
      =\begin{cases}(p^{n+1}-q^{n+1})/(p-q)&p\neq q,\\(n+1)p^n&p=q.\end{cases}
    \]
\end{minipage}

  \item \begin{minipage}[t]{\linewidth}\raggedright
\textbf{Question.} Soient $a_1,\ldots,a_n\in K$. Calculer le déterminant de la matrice $(n+1)\times(n+1)$ dont la ligne $i\leq n$ a $-a_i$ en colonne $i$, $a_i$ en colonne $i+1$, $0$ ailleurs, et dont la dernière ligne vaut $(1,\ldots,1)$.
\par\smallskip
\textbf{Réponse.} Effectuons $C_{j+1}\leftarrow C_{j+1}+C_j$ pour $j=1,\ldots,n$, dans cet ordre.
    La dernière colonne devient $(0,\ldots,0,n+1)^\top$ et les $n$ premiers coefficients diagonaux restent $-a_1,\ldots,-a_n$.
    \[
      \therefore D_n=(-1)^n(n+1)\prod_{i=1}^n a_i
    \]
\end{minipage}

  \item \begin{minipage}[t]{\linewidth}\raggedright
\textbf{Question.} Soient $a_1,\ldots,a_n\in K$. Calculer $D_n(x)=\det M(x)$ pour la matrice $(n+1)\times(n+1)$ de diagonale $x$ et de coefficients hors diagonale $m_{i,j}=a_j$ si $i>j$ et $m_{i,j}=a_{j-1}$ si $i<j$ (indices $1,\ldots,n+1$).
\par\smallskip
\textbf{Réponse.} Posons $s=x+\sum_i a_i$. On effectue $C_1\leftarrow\sum_j C_j$ : la première colonne est $s\boldsymbol{1}$.
    Par multilinéarité, on sort le facteur $s$, puis on retranche $a_jC_1$ à $C_{j+1}$.
    La matrice obtenue est triangulaire inférieure, de diagonale $1,x-a_1,\ldots,x-a_n$.
    \[
      \therefore D_n(x)=\left(x+\sum_{i=1}^n a_i\right)\prod_{i=1}^n(x-a_i)
    \]
\end{minipage}

  \item \begin{minipage}[t]{\linewidth}\raggedright
\textbf{Question.} Calculer le déterminant de la matrice $n\times n$ de diagonale $(x_1,\ldots,x_n)$, de coefficients $a$ au-dessus de la diagonale et $b$ au-dessous. Indication : ajouter $t$ à chaque coefficient et interpoler le déterminant, affine en $t$.
\par\smallskip
\textbf{Réponse.} Ajouter $t$ à chaque coefficient donne un déterminant affine $F(t)=D_n+ct$, car les termes contenant deux colonnes constantes sont nuls.
    Les matrices pour $t=-a$ et $t=-b$ sont triangulaires, donc
    $F(-a)=\prod_i(x_i-a)$ et $F(-b)=\prod_i(x_i-b)$.
    \[
      \therefore D_n=\frac{b\prod_i(x_i-a)-a\prod_i(x_i-b)}{b-a}\quad(a\neq b)
    \]
    Pour $a=b$, la matrice est diagonale plus une matrice de rang $1$ ; la multilinéarité donne
    \[
      D_n=\prod_i(x_i-a)+a\sum_{k=1}^n\prod_{i\neq k}(x_i-a)
    \]
    Cette formule reste valable lorsque certains $x_i=a$.
\end{minipage}

  \item \begin{minipage}[t]{\linewidth}\raggedright
\textbf{Question.} Soit $A\in\ZZ^{n\times n}$ avec $n$ pair, une diagonale nulle et des coefficients hors diagonale égaux à $\pm1$. Montrer que $A$ est inversible.
\par\smallskip
\textbf{Réponse.} Modulo $2$, $A\equiv J_n-I_n$. Si $(J_n-I_n)x=0$ sur $\ZZ/2\ZZ$, alors
    $x=(\sum_i x_i)\boldsymbol{1}$.
    Comme $n$ est pair, la somme des coordonnées de ce vecteur vaut $0$, donc $x=0$.
    Ainsi $\det A\equiv1\pmod2$.

    $\therefore\det A\neq0$ et $A$ est inversible sur $\QQ$ (donc sur $\RR$).
\end{minipage}

  \item \begin{minipage}[t]{\linewidth}\raggedright
\textbf{Question.} Montrer que $A\in\ZZ^{n\times n}$ a un inverse à coefficients entiers si et seulement si $\det A=\pm1$.
\par\smallskip
\textbf{Réponse.} Si $A^{-1}$ est entière, $\det A$ et $\det(A^{-1})$ sont des entiers de produit $1$, donc $\det A=\pm1$.
    Réciproquement, si $\det A=\pm1$, la formule
    $A^{-1}=\operatorname{Com}(A)^\top/\det A$ donne une matrice entière.
\end{minipage}

  \item \begin{minipage}[t]{\linewidth}\raggedright
\textbf{Question.} Soient $A,B,C,D\in K^{n\times n}$ avec $AC=CA$. Montrer que $\det\begin{pmatrix}A&B\\C&D\end{pmatrix}=\det(AD-CB)$. Indication : commencer par $A$ inversible.
\par\smallskip
\textbf{Réponse.} Si $A$ est inversible, l'élimination par blocs donne
    \[
      \det\begin{pmatrix}A&B\\C&D\end{pmatrix}
      =\det A\det(D-CA^{-1}B)=\det(AD-CB)
    \]
    La dernière égalité utilise $AC=CA$.
    Dans le cas général, on remplace $A$ par $A+tI_n$, inversible sur $K(t)$ et commutant encore avec $C$.
    L'égalité est une identité polynomiale en $t$ ; en évaluant en $0$, on obtient le résultat sur tout corps $K$.
\end{minipage}

  \item \begin{minipage}[t]{\linewidth}\raggedright
\textbf{Question.} Soient $A=\begin{pmatrix}a&b\\c&d\end{pmatrix}\in K^{2\times2}$ et $M\in K^{m\times m}$. Calculer $\det(A\otimes M)$ en fonction de $\det A,\det M$, où $A\otimes M=\begin{pmatrix}aM&bM\\cM&dM\end{pmatrix}$.
\par\smallskip
\textbf{Réponse.} Les blocs $aM$ et $cM$ commutent. D'après la question précédente,
    \[
      \det(A\otimes M)=\det(adM^2-cbM^2)
      =(ad-bc)^m\det(M)^2
      =\det(A)^m\det(M)^2
    \]
\end{minipage}

\end{enumerate}
%%%%%%%%%%%%%%%%%%%%%%%%%%%%%%%%%%%%%%%%%%%%%%%%%%%%%%
%                       End                          %
%%%%%%%%%%%%%%%%%%%%%%%%%%%%%%%%%%%%%%%%%%%%%%%%%%%%%%

\end{document}
